
% =====================================================================
\section{Findings: the order book, and the day within the day}
\label{sec:book}
% =====================================================================

\subsection{Rebuilding the paper's equations without its proxies}

\begin{table}[htbp]
\centering
\caption{Order-book anatomy of price clustering}
\label{tab:rq3}
\small
\begin{threeparttable}
\begin{tabular}{@{}lrrrrrrr@{}}
\toprule
Dependent variable & $L^{0}$ & SE & $L^{0C}$ & SE & RDepth & SE & Stock-days \\
\midrule
$M^{BLarge0}$ & 0.0029 & (0.0524) & -0.0092 & (0.0136) & 0.2565*** & (0.0222) & 9,527 \\
$\Delta\mathit{Imp}^{BLarge}$ & 3.6734 & (9.3067) & -0.7720 & (2.0219) & 17.6809* & (10.0115) & 5,366 \\
$M^{SLarge0}$ & 0.1501 & (0.0983) & 0.0117 & (0.0197) & 0.1763*** & (0.0408) & 9,537 \\
$\Delta\mathit{Imp}^{SLarge}$ & 6.0726 & (7.7144) & -2.4585 & (2.1336) & 4.1140 & (4.0574) & 5,722 \\
\bottomrule
\end{tabular}
\begin{tablenotes}[flushleft]\footnotesize
\item Analogues of Ohta's (2026) equations (1) and (2), with book-derived variables replacing the margin-trading and ownership proxies, which are not available here. $L^{0}$ is the round-price share of limit-order volume submitted within the visible book, $L^{0C}$ the ratio of cancelled to submitted volume at round prices, and RDepth the round-price share of visible ten-level resting depth. Sell-side book variables are paired with buy-initiated outcomes, since a buy market order executes against sell limit orders. All regressors are lagged one day. Stock and day fixed effects, standard errors clustered by stock and by day. The hypothesis predicts a positive coefficient on $L^{0}$ and RDepth and a negative one on $L^{0C}$. $^{*}p<0.1$, $^{**}p<0.05$, $^{***}p<0.01$.
\end{tablenotes}
\end{threeparttable}
\end{table}


Ohta explains clustering with limit-order submissions, their cancellation rate,
and two proxies for individual activity. The proxies are unavailable here, so
Table~\ref{tab:rq3} substitutes what the book itself shows: the round-price
share of submitted limit-order volume, the ratio of cancellations to submissions
at round prices, and the round-price share of standing visible depth.

The hypothesis makes three sign predictions, and one of the
three is confirmed at the 5\% level. More
round-price submissions should mean more round-price trading, so the
coefficient on $L^{S0}$ should be positive: it is 0.0041,
not statistically distinguishable from zero. Orders that are \emph{not} cancelled are the ones that get
picked off, so the coefficient on the cancellation ratio should be negative:
it is -0.0098, not statistically distinguishable from zero. And more standing round-price
inventory should mean more round-price execution: the coefficient on
$\mathit{RDepth}$ is 0.2557, significant at the 1\% level.

The two flow-based nulls need a caveat before they are read as evidence against the mechanism. The inference watches ten visible levels while 81\% of resting volume sits beyond them (Appendix~\ref{app:ladder}), so the daily $L$ and $L^{C}$ series carry heavy measurement error, and classical measurement error drives coefficients toward zero. A null here is what censoring would produce even if the mechanism held; these estimates do not adjudicate it. The depth variable below requires no inference at all, which is why it is the cleaner test.

\begin{table}[htbp]
\centering
\caption{Standing round-price depth and next-day liquidity}
\label{tab:rdepth}
\small
\begin{threeparttable}
\begin{tabular}{@{}lrrrrr@{}}
\toprule
Outcome & RDepth ask & SE & RDepth bid & SE & Stock-days \\
\midrule
Effective spread (log) & 0.1314*** & (0.0169) & 0.0433** & (0.0208) & 74,218 \\
Price impact, 60s (bp) & 1.1252*** & (0.4207) & -0.3255 & (0.4168) & 74,218 \\
Depth at best (log) & -0.0153 & (0.0231) & 0.0522* & (0.0290) & 74,218 \\
\bottomrule
\end{tabular}
\begin{tablenotes}[flushleft]\footnotesize
\item RDepth is the share of visible ten-level resting volume sitting at round prices, time-averaged over the day and lagged one day. Dynamic specification, controls and fixed effects as in Table~\ref{tab:rq2}.
\end{tablenotes}
\end{threeparttable}
\end{table}


Table~\ref{tab:rdepth} asks the more direct question --- whether standing
round-price depth predicts liquidity by itself. This is the measure with the
cleanest interpretation in the whole study, because it requires no inference at
all: it is a share of displayed volume, read off the book.

One reading discipline before celebrating that coefficient: standing depth predicting execution at the same prices is partly mechanical --- trades print where volume rests, whoever put it there and for whatever reason. The result is consistent with the stale-order account rather than diagnostic of it.

\subsection{The shape of the trading day}

Before turning to the regressions, the raw intraday profile is worth reporting on its own, because it reproduces a result from a different paper by the same author. \textcite{Ohta2006Intraday} finds that clustering on this exchange is highest at the open and decays through the session, which is what the price-resolution hypothesis predicts: uncertainty about value is greatest when trading starts, and a round number is the cheapest available substitute for an estimate.

Figure~\ref{fig:intraday} shows the same curve two decades later. The round-price share opens at 16.4\% and falls to 11.5\% in the closing bucket, a decline of 4.9 percentage points, and it does so nearly monotonically. The effective spread traces the familiar intraday pattern beside it. Neither series was tuned to produce this, and the agreement with a 2006 result computed from a different sample is a further sign that the measurement is reading the phenomenon rather than an artefact.

It also sharpens the interpretation of what follows. If clustering were simply a proxy for wide spreads, the two curves would be redundant; they are not, and the regressions below control for the spread directly.

\begin{figure}[htbp]
\centering
\includegraphics[width=0.85\linewidth]{f_intraday.pdf}
\caption{Round-price share of traded volume and the effective half-spread, by
thirty-minute bucket. The dashed horizontal line is the 10\% a uniform digit
distribution would give.}
\label{fig:intraday}
\end{figure}

\subsection{Within the day}

\begin{table}[htbp]
\centering
\caption{Clustering and liquidity within the trading day}
\label{tab:intraday}
\small
\begin{threeparttable}
\begin{tabular}{@{}lrrrr@{}}
\toprule
Outcome & Regressor & Coefficient & SE & Bucket-rows \\
\midrule
Effective spread (log) & $M^{0}_{b-1}$ & 0.0189*** & (0.0025) & 612,672 \\
Effective spread (log) & $M^{BLarge0}_{b-1}$ & 0.0242*** & (0.0053) & 252,848 \\
Price impact, 60s (bp) & $M^{0}_{b-1}$ & 0.1041 & (0.0737) & 612,670 \\
Price impact, 60s (bp) & $M^{BLarge0}_{b-1}$ & 0.0655 & (0.0508) & 252,848 \\
\bottomrule
\end{tabular}
\begin{tablenotes}[flushleft]\footnotesize
\item Thirty-minute buckets. Stock-day and bucket-of-day fixed effects, so all daily variation and the average intraday pattern are absorbed and only within-day, bucket-to-bucket variation identifies the coefficient. Standard errors clustered by stock and by date. Each specification includes the outcome's own lagged bucket and log turnover. $^{*}p<0.1$, $^{**}p<0.05$, $^{***}p<0.01$.
\end{tablenotes}
\end{threeparttable}
\end{table}


The paper asks for the relationship ``at a daily or shorter frequency''.
Table~\ref{tab:intraday} is the shorter one: thirty-minute buckets, with
stock-day fixed effects absorbing everything about the day --- the stock, the
news, the market state --- and bucket-of-day fixed effects absorbing the familiar
intraday pattern in spreads and volume. What identifies the coefficient is
purely within-day, bucket-to-bucket variation; standard errors are clustered
by stock and by date, where the dependence actually lives.

The coefficient on lagged-bucket clustering in the effective-spread regression is
0.0189, significant at the 1\% level. Coefficients here are small by construction and
should be read that way: this design deliberately throws away all the variation
the daily panel uses. This is the version of the liquidity relationship that
survives every re-specification tried in this study.

\paragraph{Order flow, conditional on clustering.}
The execution-relevant version of the impact result asks whether the \emph{same}
order-flow imbalance moves the midquote more on high-clustering days. The
interaction term is 0.0514, not statistically distinguishable from zero. A positive value would mean the book is easier to push when round-price trading is heavy; the estimate here does not resolve the question in either direction.

\paragraph{Verdict.}
Of the paper's mechanism variables, the only one confirmed is the one that requires no inference: standing round-price depth. The inferred flow shares are nulls under heavy censoring, which is uninformative rather than damning. The within-day relationship stands (0.0189, $t=8.80$) in a design where nothing slow-moving about a stock can produce it.
